\lecture{8.5}{Failure of Naive Regression}
\label{lec:8-5}

\subsection{Regression Coefficients in Single-Timepoint Settings}

Even in a point-treatment setting, a regression coefficient is not automatically
an average treatment effect. Suppose
\[
  \E(Y\given A=a,W=w)=\beta_0+\beta_1 a+\beta_2 w.
\]
Then
\[
  \E(Y\given A=1,W=w)-\E(Y\given A=0,W=w)=\beta_1
\]
for every $w$, and the standardized mean difference is also $\beta_1$:
\[
  \E\{\E(Y\given A=1,W)-\E(Y\given A=0,W)\}=\beta_1.
\]
This special equality depends on the linear model and the absence of treatment
effect modification by $W$.

\subsection{Linear Models Without Interactions}

When a linear model has no interaction between treatment and covariates, the
treatment coefficient is a constant conditional mean contrast. Under the usual
identification assumptions, it can equal the g-computation estimate of the ATE.
This is a special case, not a general rule.

\subsection{Linear Models With Interactions}

If the conditional mean includes an interaction,
\[
  \E(Y\given A=a,W=w)=\beta_0+\beta_1a+\beta_2w+\beta_3aw,
\]
then the conditional contrast is
\[
  \E(Y\given A=1,W=w)-\E(Y\given A=0,W=w)=\beta_1+\beta_3w.
\]
The marginal standardized contrast is
\[
  \E(\beta_1+\beta_3W)=\beta_1+\beta_3\E(W).
\]
Thus $\beta_1$ is the conditional contrast at $W=0$, not the average treatment
effect unless $\E(W)=0$ or $\beta_3=0$.

\subsection{Logistic Models}

For binary outcomes, logistic regression is nonlinear:
\[
  \E(Y\given A=a,W=w)=\expit(\beta_0+\beta_1a+\beta_2w).
\]
The risk difference at covariate value $w$ is
\[
  \expit(\beta_0+\beta_1+\beta_2w)-\expit(\beta_0+\beta_2w),
\]
and the marginal risk difference is
\[
  \E\left[
    \expit(\beta_0+\beta_1+\beta_2W)-\expit(\beta_0+\beta_2W)
  \right].
\]
The coefficient $\beta_1$ is a conditional log-odds ratio parameter, not a risk
difference and not a marginal ATE.

\subsection{Time-Varying Confounder Feedback}

Longitudinal settings introduce a deeper problem. Consider
\[
  L_0 \rightarrow A_0 \rightarrow L_1 \rightarrow A_1 \rightarrow Y,
\]
where $L_1$ also affects $Y$. A regression of $Y$ on $A_0,L_1,A_1$ conditions
on $L_1$. This may be useful for estimating the effect of $A_1$, because $L_1$
confounds the $A_1$--$Y$ relation. But it blocks part of the effect of $A_0$ on
$Y$ that operates through $L_1$.

\begin{figure}[H]
\centering
\begin{tikzpicture}[
  node distance=1.35cm,
  box/.style={draw, rounded corners=2pt, minimum width=1.0cm, minimum height=.72cm},
  arr/.style={-{Latex[length=2.3mm]}, thick},
  block/.style={draw=red!65!black, line width=1pt, circle, minimum size=.82cm}
]
\node[box] (A0) {$A_0$};
\node[box, right=of A0] (L1) {$L_1$};
\node[box, right=of L1] (A1) {$A_1$};
\node[box, right=of A1] (Y) {$Y$};
\draw[arr] (A0) -- (L1);
\draw[arr] (L1) -- (A1);
\draw[arr] (A1) -- (Y);
\draw[arr, bend left=23] (L1) to (Y);
\draw[arr, bend left=35] (A0) to (Y);
\node[block] at (L1) {};
\end{tikzpicture}
\caption{Conditioning on $L_1$ can block the mediated path from $A_0$ to $Y$ while controlling confounding for $A_1$.}
\label{fig:ch8-naive-blocking}
\end{figure}

\subsection{A Numerical Data-Generating Example}

Suppose
\[
  A_0\sim \operatorname{Bernoulli}(0.5),
\]
\[
  L_1\given A_0=a_0 \sim N(1+a_0,1),
\]
\[
  A_1\given L_1=l_1,A_0=a_0
  \sim \operatorname{Bernoulli}\{\expit(-1+l_1+a_0)\},
\]
and
\[
  Y\given A_1=a_1,L_1=l_1,A_0=a_0
  \sim N(1+a_1+2l_1,1).
\]
The mean outcome regression does not contain $A_0$ directly. Nevertheless,
$A_0$ has a causal effect because it shifts $L_1$, and $L_1$ affects $Y$.

Using g-computation,
\[
\begin{aligned}
  \E\{Y(1,1)\}
  &=\E\{\E(Y\given A_1=1,L_1,A_0=1)\given A_0=1\}\\
  &=\E(2+2L_1\given A_0=1)=2+2(2)=6,
\end{aligned}
\]
because $\E(L_1\given A_0=1)=2$. Similarly,
\[
  \E\{Y(1,0)\}=5,\qquad
  \E\{Y(0,1)\}=4,\qquad
  \E\{Y(0,0)\}=3.
\]
Therefore the effect of changing $A_0$ from 0 to 1 while holding $A_1=1$ is
\[
  \E\{Y(1,1)-Y(0,1)\}=2.
\]

Now fit the correctly specified conditional mean regression
\[
  \E(Y\given A_1,L_1,A_0)=1+A_1+2L_1+0\cdot A_0.
\]
The coefficient on $A_0$ is zero because the regression conditions on $L_1$.
Reading the $A_0$ coefficient as the causal effect of baseline treatment would
therefore be wrong.

\subsection{Blocking Causal Pathways by Conditioning on Intermediates}

The problem is not that regression is useless. The problem is using one
conditional regression as if it were the causal estimand. Sequential
g-computation uses regressions carefully: it conditions on $L_1$ in the final
outcome regression and then averages over the distribution of $L_1$ that would
arise under each baseline treatment. A naive regression coefficient conditions
on $L_1$ and stops there.

\begin{keyidea}
In longitudinal studies, the same covariate can be a confounder for later
treatment and a mediator for earlier treatment. Ordinary regression adjustment
cannot generally handle this feedback by reading off coefficients.
\end{keyidea}
